\section{Occupation above rank and finite-group rigidity}\label{sec:rigidity51}
We retain the hypotheses of Section~\ref{sec:lifts51}. For an integer
$V\ge D+1$ put
\[
 \mathscr P_{V,\rho}=\{S:C_\rho\subset S\subset[-1,1]^D,
                  \ S\text{ is a polytope},\ f_0(S)\le V\}.
\]
With $I\in\calA$ still assumed, whenever this class is nonempty, define
\begin{equation}\label{eq:gamma51}
 \gamma_{\calA}(V,\rho)=
 \min_{S\in\mathscr P_{V,\rho}}
 \frac{\operatorname{vol}(\conv\bigcup_{a\in\calA}U_aS)}
      {\operatorname{vol}(S)}.
\end{equation}
This is a finite-dimensional variational constant, not a claimed
closed formula for optimal stochastic width. We use
$V_K=\binom{K}{\lfloor K/2\rfloor}$; the sharper rank-dependent
bound of Theorem~\ref{thm:section51} can replace it whenever the ranks
are known.

\begin{lemma}[Compactness and strict expansion]\label{lem:compact51}
The class $\mathscr P_{V,\rho}$ is compact in Hausdorff distance.
The minimum in \eqref{eq:gamma51} is attained and is at least one.
It equals one if and only if the class contains a polytope invariant
under every $U_a$. Such a polytope gives a faithful permutation action
of $G=\langle U_a:a\in\calA\rangle$ on at most $V$ vertices.
In particular, if $G$ is infinite, $\gamma_{\calA}(V,\rho)>1$.
\end{lemma}
\begin{proof}
Pad each vertex list to length $V$ by repetitions. Subsequences of
these lists converge in the compact cube; their convex hulls converge
to a polytope with at most $V$ vertices. Containment of $C_\rho$
passes to the limit, so no dimension is lost. This proves compactness.
Finite unions followed by convex hull commute with Hausdorff limits.
Volume is continuous here: every body contains the ball of radius
$\rho/\sqrt D$, and a Hausdorff perturbation by $h$ is trapped
between dilates with factors tending to one as $h\to0$.
The denominators are bounded below by
$\operatorname{vol}(C_\rho)=2^D\rho^D/D!$. Thus the minimum is attained.

Since the alphabet contains $I$, each hull in the numerator contains
$S$. Equality of volumes forces equality of the two convex bodies.
Indeed a proper inclusion of full-dimensional compact convex bodies
contains, by separation and convexity with an interior ball, an extra
set of positive volume. Hence $U_aS\subset S$. Orthogonality and the
same strictness argument give $U_aS=S$. The converse is immediate.
Each map permutes the vertices. Any group element fixing them all
fixes their affine hull, which is $\R^D$, and is the identity. This
proves faithfulness and the extremely coarse factorial bound $|G|\le V!$.
\end{proof}

\begin{theorem}[All-width occupation budget]\label{thm:occupation51}
Fix $K\ge D+1$. In an arbitrary exact realization let $m\ge1$ be
the number of command cuts with available width at most $K$.
If $\mathscr P_{V_K,\rho}$ is empty, no such cut exists. Otherwise
\begin{equation}\label{eq:occupation51}
 \gamma_{\calA}(V_K,\rho)^{m-1}\le D!\rho^{-D}.
\end{equation}
All intervening registers may have arbitrary finite width. If $G$ is
infinite, this bounds the occupation of every fixed finite width by a
constant independent of the horizon.
\end{theorem}
\begin{proof}
Order the selected cuts $s_1<\cdots<s_m$ and apply
Theorem~\ref{thm:section51}. Since central binomial coefficients
increase with their upper index, each physical polytope $S_{s_i}$
has at most $V_K$ vertices. For every interval between selected cuts,
a word with one specified letter $a$ and identities elsewhere has
product $U_a$. Consequently
\[
 \conv\bigcup_{a\in\calA}U_aS_{s_i}\subset S_{s_{i+1}}.
\]
Equation~\eqref{eq:gamma51} multiplies volume by at least
$\gamma_{\calA}(V_K,\rho)$ at each interval. The first volume is at
least $2^D\rho^D/D!$ and the last is at most $2^D$.
Their ratio proves \eqref{eq:occupation51}. Strictness for an infinite
group follows from Lemma~\ref{lem:compact51}.
\end{proof}

\begin{theorem}[Finite-group criterion and stationarization]
\label{thm:finitegroup51}
The following statements are equivalent:
\begin{enumerate}
\item The peaks $W_{N,0}$ are bounded as $N$ varies, allowing a different
clocked machine and different real rows at every horizon.
\item There is one finite-state exact machine whose command updates
are permutations and which works at every horizon.
\item The orthogonal group $G$ generated by the alphabet is finite.
\end{enumerate}
More precisely, a uniform clocked bound $K$ gives an invariant physical
polytope with at most $V_K$ vertices, a stationary permutation machine
with at most $V_K$ labels, and $|G|\le V_K!$. Conversely a finite
group admits such a machine with at most $2D|G|$ labels. If $G$ is
infinite, $W_{N,0}\to\infty$.
\end{theorem}
\begin{proof}
Suppose the bound $K$ in (1) holds. The class
$\mathscr P_{V_K,\rho}$ is nonempty. Applying
Theorem~\ref{thm:occupation51} to all $N+1$ command cuts gives
$\gamma_{\calA}(V_K,\rho)^N\le D!\rho^{-D}$ for every $N$.
It follows that $\gamma_{\calA}(V_K,\rho)=1$.
Lemma~\ref{lem:compact51} supplies an invariant polytope $S$ with at
most $V_K$ vertices. Initialize $\rho x$ barycentrically in its
vertices, update by the induced vertex permutations, and decode a
coordinate by its vertex value. All values lie in $[-1,1]$.
This proves (2), including its quantitative bound, and the faithful
action proves (3).

If (3) holds, store a point of the finite orbit
$\{gx:g\in G,x\in X\}$, initialize at $x$, update by $U_a$, and
decode with mean $\rho e_j^Tgx$. This uses at most $2D|G|$ labels,
is exact, and is independent of both time and horizon. Thus (3)
implies (2), and (2) implies (1). The held binary answer requires its
two labels at its own cut and does not increase these peaks.
Finally $W_{N,0}$ is nondecreasing: a machine for $N+1$ can answer
at length $N$ by incorporating one final identity update into its
decoder, explicitly replacing $d_j$ by $T_{N+1,I}d_j$.
If $G$ is infinite the sequence is unbounded by (1)--(3),
so it tends to infinity.
\end{proof}

The stationary bound need not equal $K$: passage from a section to
its projected vertices can enlarge the register. Thus the theorem
does not assert that optimal hidden width equals optimal physical
vertex count. Nor does it infer a sharp rate from the qualitative
strictness of \eqref{eq:gamma51}.

\begin{corollary}[A rational noncommuting family in fixed dimension]
\label{cor:rational51}
In dimension two take $X=\{\pm e_1,\pm e_2\}$ and
\[
 R=\begin{pmatrix}3/5&-4/5\\4/5&3/5\end{pmatrix},\qquad
 F=\begin{pmatrix}1&0\\0&-1\end{pmatrix},\qquad
 \calA=\{I,R,F\}.
\]
For every fixed $0<\rho<1$, $W_{N,0}\to\infty$, and for each
fixed $K$ the occupation bound \eqref{eq:occupation51} has a multiplier
strictly greater than one. More generally the same conclusions hold
for $\{I,R_\theta,F\}$ whenever $\theta/\pi$ is irrational.
\end{corollary}
\begin{proof}
One has $FRF=R^{-1}\ne R$, so the commands do not commute. If
$\zeta=(3+4i)/5$ were a root of unity, then
$\zeta+\zeta^{-1}=6/5$ would be a rational algebraic integer (the trace, or sum of the two
eigenvalues),
hence an integer, a contradiction. Here the last elementary fact
follows by clearing denominators in a monic integer equation for a
rational number written in lowest terms. Thus $R$ has infinite order.
Apply Theorems~\ref{thm:occupation51} and~\ref{thm:finitegroup51}.
An irrational-angle rotation has infinite order by its definition,
and the reflection relation is unchanged.
\end{proof}

Dimension, alphabet and query set in this corollary are fixed while
$N$ tends to infinity. It is distinct from the finite tensor embedding
of Section~\ref{sec:orthogonal50}, whose dimension grows with the
word table. The variational bound also differs from the sharper
arithmetic estimates available for particular commuting rotations in
the later sections.

\subsection{An antipodal budget for one surplus label}
A further consequence of the reachability rank is available even for
the finite group $\{I,-I\}$. For $D\ge3$ let
\[
 \mathscr H_{D,\rho}=\{S:C_\rho\subset S\subset[-1,1]^D,
       \ f_0(S)\le D+2\ \hbox{or}\ f_{D-1}(S)\le D+2\}.
\]
Here all members are full-dimensional polytopes. If this class is
nonempty put
\begin{equation}\label{eq:antipodal51}
 \beta_D(\rho)=\min_{S\in\mathscr H_{D,\rho}}
       \frac{\operatorname{vol}(\conv(S\cup(-S)))}
            {\operatorname{vol}(S)}.
\end{equation}

\begin{theorem}[One-surplus occupation and a six-state frontier]
\label{thm:surplus-budget51}
For $D\ge3$ the minimum \eqref{eq:antipodal51}, whenever defined,
is attained and satisfies $\beta_D(\rho)>1$.
If $I,-I\in\calA$, and an exact machine has $m\ge1$ command cuts
of width at most $D+2$, then this class is nonempty and
\begin{equation}\label{eq:surplus-budget51}
 \beta_D(\rho)^{m-1}\le D!\rho^{-D}.
\end{equation}
In dimension three, for any signed-permutation alphabet containing
$I,-I$, one consequently has
\begin{equation}\label{eq:six51}
 W_{N,0}=6\qquad\hbox{whenever}\qquad
 \beta_3(\rho)^N>6\rho^{-3}.
\end{equation}
Thus six labels are exactly optimal for all sufficiently large
horizons in this class, including noncommuting alphabets.
\end{theorem}
\begin{proof}
The subclass with at most $D+2$ vertices is compact by
Lemma~\ref{lem:compact51}. For the facet subclass take polar bodies
about zero, which lies in every interior because $C_\rho$ is contained. They satisfy
$C_1\subset S^\circ\subset\rho^{-1}[-1,1]^D$ and have at most
$D+2$ vertices. Vertex compactness and continuity of polarity between
these fixed inner and outer bounds prove compactness of that
subclass as well. The map $S\mapsto\conv(S\cup(-S))$ is continuous
in Hausdorff distance, and its volume is continuous under these fixed
inner and outer bounds. Hence their union is compact and the minimum is
attained.

A centrally symmetric full-dimensional polytope has at least $2D$
vertices: its antipodal vertex pairs must linearly span $\R^D$.
It also has at least $2D$ facets, either by polarity or by applying
the same argument to paired facet normals. Since $D+2<2D$, no member
of $\mathscr H_{D,\rho}$ equals its negative. Strict volume
monotonicity therefore makes every ratio in
\eqref{eq:antipodal51} greater than one; compactness makes the
minimum strictly greater than one.

At each selected cut, rank excludes fewer than $D+1$ labels.
Corollary~\ref{cor:surplus51} puts its physical polytope in
$\mathscr H_{D,\rho}$, whether its width is $D+1$ or $D+2$.
Identity padding and one antipodal command imply that the next
selected polytope contains both signs of the previous one.
Iteration and the same endpoint volume bounds as before prove
\eqref{eq:surplus-budget51}.

For $D=3$ the class is nonempty for every $0<\rho<1$: the tetrahedron
with vertices $(1,1,1)$, $(1,-1,-1)$, $(-1,1,-1)$, $(-1,-1,1)$
lies in the cube and contains $C_1$, since each $\pm e_i$ is an
edge midpoint. A peak of at most five would violate
\eqref{eq:surplus-budget51} under the condition in \eqref{eq:six51}.
Six labels suffice by storing $\{\pm e_1,\pm e_2,\pm e_3\}$,
permuting under commands, and decoding with mean $\rho x_j$.
This proves the exact equality under a sufficient variational horizon
condition. The effective lower bound and a numerical sufficient horizon
are supplied in Theorem~\ref{thm:beta52}; a sharp first horizon is not
asserted by either result.
\end{proof}
